Teams of robots often have to coordinate when each member sees only part of the world. Learned communication protocols, either continuous message vectors trained by backpropagation~\cite{sukhbaatar2016learning,foerster2016learning,das2018tarmac} or discrete symbols~\cite{lazaridou2016multi,havrylov2017emergence,mordatch2017emergence}, solve such tasks but tend to be private conventions: they need not resemble natural language~\cite{kottur2017natural}, and success inside a team can overstate what the messages mean~\cite{lowe2019pitfalls}. A fixed, human-readable language would be legible and would not depend on the training partner, but it constrains what can be said.

We ask a narrow question at a scale that fits on a laptop CPU: in a controlled rendezvous task with partial observability, how do (a) no communication, (b) a learned continuous vector, (c) a learned discrete token channel and (d) a fixed templated language compare on final success, sample efficiency and cross-play with a partner trained separately? The templated language is a small-scale \emph{proxy} for natural language: a two-word symbolic template with a fixed lexicon. It captures legibility and partner independence, but not ambiguity, open vocabulary or pragmatics, and we make no claim about natural language itself.

Contributions are modest: a fully reproducible comparison (code, run registry, 5 seeds) with channels of nominally matched capacity for the two symbolic channels (the continuous vector is not matched; it carries far more information), a cross-play protocol, and capacity and noise sweeps. Results are mixed for the proxy language: it is robust to partner swap but not the most successful or most sample-efficient system here.
